\chapter{Álgebra partida y factorización exacta del carácter}
\label{chap:factorizacion-split}
\index{álgebra partida}\index{carácter!factorización partida}
\index{cono nulo}\index{canales $q_\pm$}

La firma formal de una extensión y los canales angulares proceden de ramas
ya construidas. En esta sección ambas confluyen en una factorización
determinantal que conserva, por separado, la escala radial y la orientación
relativa. El resultado pertenece al estrato algebraico: no identifica por sí
solo los ideales nulos con fotones físicos ni el interior positivo con una
partícula material.

\section{Coordenadas angulares conjugadas y retícula espectral bidireccional}
\label{sec:reticula-angular-split}

La cadena causal anterior fija
\[
 A_{\rm rad}=\frac{\pi A}{180},\qquad
 \Cstarrad=\frac{\pi\Cstar}{180},
\]
y, sólo después de obtener la coordenada directa \(A\) y la coordenada
conjugada única \(C^*\),
\[
 q_+=\exp[-(A_{\rm rad}+\Cstarrad)],
 \qquad
 q_-=\exp[-(A_{\rm rad}-\Cstarrad)].
\]
Para una firma angular \((n_A,s_C)\), introducimos los dos números de
enrollamiento
\[
 W_+=6n_A+s_C,
 \qquad
 W_-=6n_A-s_C.
\]

\begin{theorem}[Retícula angular de índice doce]
\label{thm:reticula-two-way}
La aplicación
\[
 \mathcal W:\ZZ^2\to\ZZ^2,
 \qquad
 \binom{n_A}{s_C}\longmapsto
 \binom{W_+}{W_-}
 =\begin{pmatrix}6&1\\6&-1\end{pmatrix}
 \binom{n_A}{s_C}
\]
es inyectiva y su imagen tiene índice doce. Sus divisores invariantes son
\((1,12)\), y
\[
 \operatorname{im}\mathcal W
 =\{(u,v)\in\ZZ^2:u+v\equiv0\pmod{12}\}.
\]
Además,
\[
 \boxed{
 \exp\!\left(n_AA_{\rm rad}+\frac{s_C}{6}\Cstarrad\right)
 =\left(q_+^{-W_+}q_-^{-W_-}\right)^{1/12}.}
\]
\end{theorem}

\begin{proof}
La matriz tiene determinante \(-12\). Como el máximo común divisor de sus
entradas es uno, su forma normal de Smith es
\(\operatorname{diag}(1,12)\). Si
\((u,v)=(6n+s,6n-s)\), entonces \(u+v=12n\). Recíprocamente, si
\(u+v\) es múltiplo de doce, entonces
\[
 n=\frac{u+v}{12},
 \qquad
 s=\frac{u-v}{2}
\]
son enteros y recuperan \((u,v)\). Finalmente,
\[
 \frac{W_+(A_{\rm rad}+\Cstarrad)
       +W_-(A_{\rm rad}-\Cstarrad)}{12}
 =n_AA_{\rm rad}+\frac{s_C}{6}\Cstarrad,
\]
lo que prueba la identidad al exponenciar.
\end{proof}

La involución \(s_C\mapsto-s_C\) intercambia
\(W_+\leftrightarrow W_-\) y \(q_+\leftrightarrow q_-\). La raíz
duodécima es el índice integral de la imagen de \(\mathcal W\), no un
redondeo. La coordenada \(C^*\) actúa como componente antisimétrica de la retícula;
no es una corrección decimal de una masa ni una salida de
\(\Gamma_{6\to8}\).

\section{Álgebra real partida}
\label{sec:algebra-real-partida}

\begin{definition}[Álgebra real partida]
Sea
\[
 \mathcal A_{\rm sp}
 :=\RR[\mathsf R]/(\mathsf R^2-1).
\]
La conjugación y la norma se definen, respectivamente, por
\[
 \overline{a+b\mathsf R}=a-b\mathsf R,
 \qquad
 N_{\rm sp}(z)=z\overline z=a^2-b^2.
\]
\end{definition}

Los idempotentes conjugados
\[
 P_\pm=\frac{1\pm\mathsf R}{2}
\]
satisfacen
\[
 P_\pm^2=P_\pm,\qquad
 P_+P_-=0,\qquad
 P_++P_-=1,\qquad
 \overline{P_+}=P_-.
\]

\begin{proposition}[Descomposición espectral partida]
\label{prop:descomposicion-espectral-partida}
Todo \(z\in\mathcal A_{\rm sp}\) admite una descomposición única
\[
 z=r_+P_++r_-P_-,
 \qquad r_\pm\in\RR,
\]
y su norma es
\[
 \boxed{N_{\rm sp}(z)=r_+r_-.}
\]
Cada ideal principal \(\mathcal A_{\rm sp}P_\pm=\RR P_\pm\) es una recta
nula; en particular,
\[
 N_{\rm sp}(\lambda P_\pm)=0.
\]
\end{proposition}

\begin{proof}
Si \(z=a+b\mathsf R\), la elección \(r_+=a+b\), \(r_-=a-b\) proporciona la
descomposición, cuya unicidad se sigue de la independencia de
\(P_+\) y \(P_-\). La conjugación intercambia los dos coeficientes:
\(\overline z=r_-P_++r_+P_-\). La ortogonalidad de los idempotentes da
\[
 z\overline z=(r_+r_-)(P_++P_-)=r_+r_-.
\]
La afirmación sobre cada ideal se obtiene anulando uno de los dos
coeficientes.
\end{proof}

\begin{definition}[Cono causal partido]
\label{def:cono-causal-partido}
El cono positivo, su frontera nula y su interior son
\[
 \mathcal C_{\rm sp}^{+}
 =\{r_+P_++r_-P_-:r_+,r_-\geq0\},
\]
\[
 \partial_{\rm null}\mathcal C_{\rm sp}^{+}
 =\RR_{\geq0}P_+\cup\RR_{\geq0}P_-,
\qquad
 \operatorname{int}\mathcal C_{\rm sp}^{+}
 =\{r_+P_++r_-P_-:r_+r_->0\}.
\]
\end{definition}

La frontera posee una sola dirección activa y norma cero. El interior exige
dos direcciones simultáneas y norma positiva. Esta estratificación es exacta;
su realización física exige mapas que preserven la norma, el transporte y la
dinámica pertinentes.

\section{Transportador relativo y condiciones de realización representacional}
\label{sec:transportador-relativo-split-md}

El bloque histórico \(B_0=7I+2\mathsf R\) tiene espectro \(\{9,5\}\), y el
bloque angular \(B_1=AI+\Cstar\mathsf R\) tiene espectro
\(\{A+\Cstar,A-\Cstar\}\). El
\cref{thm:no-go-entrelazador-bloques-split} demuestra que, con los valores internos
vigentes, todo supuesto entrelazador lineal
\(\Phi B_0=B_1\Phi\) es nulo. El objeto activo es el transportador
multiplicativo relativo ya construido, no una equivalencia entre
representaciones de espectros disjuntos.

\section{Factorización partida del carácter angular}
\label{sec:factorizacion-split-angular}

Sea
\[
 \sigma=\frac{s_C}{6},\qquad
 w_\pm=\frac{n_A\pm\sigma}{2}=\frac{W_\pm}{12},
 \qquad
 r_\pm=q_\pm^{-w_\pm},
\]
y definamos
\[
 Z_v=r_+P_++r_-P_-.
\]

\begin{theorem}[Factorización partida del carácter angular]
\label{thm:split-character}
Para toda firma angular \(v=(n_A,s_C)\in\ZZ^2\),
\[
 \boxed{
 N_{\rm sp}(Z_v)
 =\det Z_v
 =\exp\!\left(
 n_AA_{\rm rad}+\frac{s_C}{6}\Cstarrad
 \right).}
\]
\end{theorem}

\begin{proof}
Por la definición de los canales,
\[
 \log r_+
 =\frac{n_A+\sigma}{2}(A_{\rm rad}+\Cstarrad),
 \qquad
 \log r_-
 =\frac{n_A-\sigma}{2}(A_{\rm rad}-\Cstarrad).
\]
Al sumar,
\[
 \log(r_+r_-)
 =n_AA_{\rm rad}+\sigma\Cstarrad.
\]
La \cref{prop:descomposicion-espectral-partida} da
\(N_{\rm sp}(Z_v)=r_+r_-\), y la representación diagonal en la base
\((P_+,P_-)\) da la misma expresión para el determinante.
\end{proof}

La componente que no ve el determinante conserva la orientación relativa:
\[
 \frac12\log\frac{r_+}{r_-}
 =\frac12\left(
 n_A\Cstarrad+\frac{s_C}{6}A_{\rm rad}
 \right),
\]
\[
 Z_v=
 \exp\!\left[
 \frac12\left(n_AA_{\rm rad}+\frac{s_C}{6}\Cstarrad\right)
 \right]
 \exp\!\left[
 \frac12\left(n_A\Cstarrad+\frac{s_C}{6}A_{\rm rad}\right)
 \mathsf R
 \right].
\]
El primer factor fija la escala radial; el segundo registra la orientación
relativa. Son dos lecturas coordinadas de un mismo elemento.

\section{Incorporación central de las torres}
\label{sec:factorizacion-split-torres}

Una vez construidos globalmente \(\Dfour\), \(s_{120}\) y
\(\zCorr=\zeta_{270}^{\rm corr}\), sea
\[
 t(v)=
 k_{\Dfour}\Dfour
 +\nu_{120}s_{120}
 +\nu_{270}\zCorr,
 \qquad
 \widetilde Z_v=e^{t(v)/2}Z_v.
\]

\begin{corollary}[Carácter completo como norma partida]
\label{cor:split-character-torres}
Fijados los bloques globales anteriores,
\[
 \boxed{
 N_{\rm sp}(\widetilde Z_v)
 =\exp\!\left(
 n_AA_{\rm rad}+\frac{s_C}{6}\Cstarrad
 +k_{\Dfour}\Dfour+\nu_{120}s_{120}
 +\nu_{270}\zCorr
 \right).}
\]
\end{corollary}

\begin{proof}
En la representación bidimensional, la multiplicación por el escalar central
\(e^{t(v)/2}\) multiplica el determinante por \(e^{t(v)}\). Se aplica el
\cref{thm:split-character}.
\end{proof}

El corolario construye la realización partida declarada del carácter activo;
no clasifica todas las realizaciones posibles ni transforma las torres en
pruebas independientes. Sus hipótesis son los canales y los bloques globales
ya fijados, y su control negativo es inmediato: omitir el factor central
\(1/2\) duplica el exponente de torre; activar simultáneamente ambos
idempotentes con coeficientes positivos abandona la frontera nula.
